\subsection{Safe Navigation in Unstructured Environment}

\textbf{Experimental Setup}
We evaluate \textbf{Q2}-\textbf{Q3} using a high-fidelity vehicle dynamics simulator~-~IPG CarMaker. For all experiments, navigation is performed using only local traversability maps generated from LiDAR point clouds, with no prior knowledge or global map of the environments. For each method, navigation is performed $15$ times.

A realistic off-road environment is designed with large bumps and randomly patterned rough terrain (see Fig.~\ref{fig:driving})to conduct navigation with traversability maps produced by various methods~(\textbf{Q2}). Traversability maps are produced from our method with and without online adaptation, an elevation-map based method~(\textit{Elevation Based})~\cite{fankhauser2014robot,miki2022elevation}, and a slope-based method~(\textit{Slope Based})~\cite{sock2016probabilistic,kim2017traversable}.

\begin{table}[b]
\vspace*{-0.1in}
\centering
\small
\renewcommand {\arraystretch}{1.15}
\caption{The average and maximum motions of the vehicle across all trials.}
\resizebox{1.0\linewidth}{!}{
\label{tab:results}
    \begin{tabular}{cccccccccc}
        \toprule
        \multirow{2}{*}{\textbf{Method}} & \multirow{2}{*}{\textbf{Success Rate}} & \multicolumn{2}{c}{\bf{Vertical Vel. [m/s]}} & \multicolumn{2}{c}{\bf{Vertical Acc. [m/s$^\text{2}$]}} & 
        \multicolumn{2}{c}{\bf{Roll Acc. [rad/s$^\text{2}$]}} &
        \multicolumn{2}{c}{\bf{Pitch Acc. [rad/s$^\text{2}$]}} \\
        \cmidrule(l{4pt}r{4pt}){3-4} \cmidrule(l{4pt}r{4pt}){5-6} \cmidrule(l{4pt}r{4pt}){7-8} \cmidrule(l{4pt}r{4pt}){9-10} & & \textbf{Mean} & \textbf{Max} & \textbf{Mean} & \textbf{Max} & \textbf{Mean} & \textbf{Max} & \textbf{Mean} & \textbf{Max}\\
        \midrule
        Elevation Based & 7 / 15 & 0.152 & 2.766 & 1.254 & 40.108 & 1.287 & 54.783 & 1.177 & 34.029\\
        Slope Based & 6 / 15 & 0.128 & 1.846 & 0.999 & 30.481 & 1.018 & 38.910 & 0.953 & 26.822\\
        Ours(w.o. Adaptation) & 14 / 15 & 0.109 & 1.690 &  0.868 & 27.948 & 0.899 & 36.022 & 0.878 & 27.046 \\
        \rowcolor[HTML]{C0C0C0}Ours(w. Adaptation) & 15 / 15 & 0.106 & 1.591 & 0.854 & 25.686 & 0.892 & 32.217 & 0.849 & 22.079\\
        \bottomrule
    \end{tabular}
    }
\label{table:safe}
\end{table}

\begin{figure}[ht]
\centering
\includegraphics[width=0.9\linewidth]{Figures_Final/driving_compressed.pdf}
\caption{The off-road race track designed in the IPG CarMaker simulator for navigation experiments. The vehicle trajectory is displayed.}
\label{fig:driving}
\vspace*{-0.25in}
\end{figure}

An additional off-road environment is designed with several unknown types of bumps (see Fig.~\ref{fig:adaptation}), to validate that our model is capable of online-adapting to unknown terrains~(\textbf{Q3}), which eventually results in safe navigation. The vehicle begins to adapt after experiencing three unseen types of bumps. This allows the vehicle to effectively optimize trajectories that minimize perturbation when re-encountering the bumps. 


Based on our prior work \cite{kim_smooth_2022}, we formulate a simple running cost function $q(\vx_t)$ of the controller, which is designed to minimize the traversability cost of a trajectory while maintaining the $30$ km/h target speed:
\begin{equation}
    q(\vx_t) = \alpha_1{\text{Track}(\vx_t)} + \alpha_2{\text{Stable}(\vx_t)} + \alpha_3{\text{Speed}(\vx_t)},
\end{equation} where $\text{Track}(\vx_t)$ and $\text{Stable}(\vx_t)$ are assigned based on the local traversability map. $\text{Track}(\vx_t)$ imposes a significant penalty based on (Sec.~\ref{sec:uncertainty}) and $\text{Stable}(\vx_t)$ is based on the traversability cost map generated by our method (Sec.~\ref{sec:costmap}).



























